% GENERATED FILE. Edit canonical lessons, then run npm run generate:books.
% canonical-source-hash: fnv1a64:7695eb0725ce8bad
% canonical-lessons: JA-C410-jimoto, JA-C410-shusshin, JA-C410-hatake, JA-C410-nohara, JA-C410-shigeru

\chapter{Where I Come From}
\label{ch:ja-where-i-come-from}
\hlchaptermodality{410}

\begin{chapteropening}
\textbf{By the end of this chapter:} \emph{I can say where I come from and describe the land there.}

Complete the last lesson of chapter 410: I can say where I come from and describe the land there.
\end{chapteropening}

\section[\texorpdfstring{\ja{じもと} (jimoto)}{じもと (jimoto)}]{\ja{じもと} (jimoto) — one's home area; the local area}

\label{lesson:JA-C410-jimoto}



\begin{quote}
Before the new one: say the Japanese for a boundary, then the Japanese for a dialect.
\end{quote}

\subsection*{You'll want to know: \ja{じもと}}
\textbf{\ja{じもと}} — \emph{jimoto} — \textquotedblleft{}one's home area; the local area\textquotedblright{}.

\begin{quote}
\textbf{\ja{じもとの} \ja{まつりは} \ja{なつです。}} — \emph{jimoto no matsuri wa natsu desu} — \textquotedblleft{}The festival back home is in summer.\textquotedblright{}
\end{quote}

\begin{cousinweb}
\textbf{\ja{じもと}} (\emph{jimoto}) is built from both layers of the vocabulary: the Sino-Japanese \emph{ji}, \textquotedblleft{}ground\textquotedblright{} (Middle Chinese \emph{dijH}), an on'yomi borrowed from Chinese, and the native Japanese (wago) \emph{moto}, \textquotedblleft{}origin\textquotedblright{}, inherited rather than borrowed (JMdict; KANJIDIC2; Middle Chinese in Baxter's transcription).
\end{cousinweb}

\begin{grammarlens}[title={What you've built}]
One more word for getting around.
\end{grammarlens}

\subsection*{Your turn}
\begin{itemize}
\raggedright
  \item \emph{Say it:} \emph{jimoto}
  \item \emph{Say it:} \emph{jimoto}, once more
  \item \emph{Say it:} \emph{jimoto no matsuri wa natsu desu} — The festival back home is in summer.
  \item \emph{Recall:} say the Japanese for a boundary, then the Japanese for a dialect, then say \emph{jimoto} again
\end{itemize}

\begin{tcolorbox}[breakable,colback=teal!4,colframe=teal!35!black,title={Before you move on}]
What does \ja{じもと} mean? (\textquotedblleft{}One's home area; the local area\textquotedblright{}.) Say it once more.
\end{tcolorbox}
\section[\texorpdfstring{\ja{しゅっしん} (shusshin)}{しゅっしん (shusshin)}]{\ja{しゅっしん} (shusshin) — one's place of origin}

\label{lesson:JA-C410-shusshin}



\begin{quote}
Before the new one: say the Japanese for a dialect, then the Japanese for one's home area.
\end{quote}

\subsection*{You'll want to know: \ja{しゅっしん}}
\textbf{\ja{しゅっしん}} — \emph{shusshin} — \textquotedblleft{}one's place of origin\textquotedblright{}.

\begin{quote}
\textbf{\ja{せんせいは} \ja{どこの} \ja{しゅっしんですか。}} — \emph{sensei wa doko no shusshin desu ka} — \textquotedblleft{}Where are you from, teacher?\textquotedblright{}
\end{quote}

\begin{cousinweb}
\textbf{\ja{しゅっしん}} (\emph{shusshin}) is Sino-Japanese (kango): two Chinese characters read with their on'yomi, readings borrowed from Chinese. They are \emph{shutsu}, \textquotedblleft{}exit\textquotedblright{} (Middle Chinese \emph{tsyhwit}) and \emph{shin}, \textquotedblleft{}somebody\textquotedblright{} (\emph{syin}), joined as \emph{shusshin} (KANJIDIC2; Middle Chinese in Baxter's transcription).
\end{cousinweb}

\begin{grammarlens}[title={What you've built}]
One more word for getting around.
\end{grammarlens}

\subsection*{Your turn}
\begin{itemize}
\raggedright
  \item \emph{Say it:} \emph{shusshin}
  \item \emph{Say it:} \emph{shusshin}, once more
  \item \emph{Say it:} \emph{sensei wa doko no shusshin desu ka} — Where are you from, teacher?
  \item \emph{Recall:} say the Japanese for a dialect, then the Japanese for one's home area, then say \emph{shusshin} again
\end{itemize}

\begin{tcolorbox}[breakable,colback=teal!4,colframe=teal!35!black,title={Before you move on}]
What does \ja{しゅっしん} mean? (\textquotedblleft{}One's place of origin\textquotedblright{}.) Say it once more.
\end{tcolorbox}
\section[\texorpdfstring{\ja{はたけ} (hatake)}{はたけ (hatake)}]{\ja{はたけ} (hatake) — a field (for crops)}

\label{lesson:JA-C410-hatake}



\begin{quote}
Before the new one: say the Japanese for one's home area, then the Japanese for one's place of origin.
\end{quote}

\subsection*{You'll want to know: \ja{はたけ}}
\textbf{\ja{はたけ}} — \emph{hatake} — \textquotedblleft{}a field (for crops)\textquotedblright{}.

\begin{quote}
\textbf{\ja{そふは} \ja{はたけで} \ja{やさいを} \ja{そだてています。}} — \emph{sofu wa hatake de yasai o sodatete imasu} — \textquotedblleft{}My grandfather grows vegetables in his field.\textquotedblright{}
\end{quote}

\begin{cousinweb}
\textbf{\ja{はたけ}} (\emph{hatake}) is a native Japanese word (wago), inherited rather than borrowed from Chinese; its kanji was attached later as a kun'yomi (JMdict).
\end{cousinweb}

\begin{grammarlens}[title={What you've built}]
One more word for getting around.
\end{grammarlens}

\subsection*{Your turn}
\begin{itemize}
\raggedright
  \item \emph{Say it:} \emph{hatake}
  \item \emph{Say it:} \emph{hatake}, once more
  \item \emph{Say it:} \emph{sofu wa hatake de yasai o sodatete imasu} — My grandfather grows vegetables in his field.
  \item \emph{Recall:} say the Japanese for one's home area, then the Japanese for one's place of origin, then say \emph{hatake} again
\end{itemize}

\begin{tcolorbox}[breakable,colback=teal!4,colframe=teal!35!black,title={Before you move on}]
What does \ja{はたけ} mean? (\textquotedblleft{}A field (for crops)\textquotedblright{}.) Say it once more.
\end{tcolorbox}
\section[\texorpdfstring{\ja{のはら} (nohara)}{のはら (nohara)}]{\ja{のはら} (nohara) — an open field; a meadow}

\label{lesson:JA-C410-nohara}



\begin{quote}
Before the new one: say the Japanese for one's place of origin, then the Japanese for a field (for crops).
\end{quote}

\subsection*{You'll want to know: \ja{のはら}}
\textbf{\ja{のはら}} — \emph{nohara} — \textquotedblleft{}an open field; a meadow\textquotedblright{}.

\begin{quote}
\textbf{\ja{こどもが} \ja{のはらを} \ja{はしっています。}} — \emph{kodomo ga nohara o hashitte imasu} — \textquotedblleft{}The children are running across the meadow.\textquotedblright{}
\end{quote}

\begin{cousinweb}
\textbf{\ja{のはら}} (\emph{nohara}) is a native Japanese word (wago), inherited rather than borrowed from Chinese, built from \emph{no}, \textquotedblleft{}field\textquotedblright{} and \emph{hara}, \textquotedblleft{}plain\textquotedblright{}; its kanji were attached later as kun'yomi (JMdict).
\end{cousinweb}

\begin{grammarlens}[title={What you've built}]
One more word for getting around.
\end{grammarlens}

\subsection*{Your turn}
\begin{itemize}
\raggedright
  \item \emph{Say it:} \emph{nohara}
  \item \emph{Say it:} \emph{nohara}, once more
  \item \emph{Say it:} \emph{kodomo ga nohara o hashitte imasu} — The children are running across the meadow.
  \item \emph{Recall:} say the Japanese for one's place of origin, then the Japanese for a field (for crops), then say \emph{nohara} again
\end{itemize}

\begin{tcolorbox}[breakable,colback=teal!4,colframe=teal!35!black,title={Before you move on}]
What does \ja{のはら} mean? (\textquotedblleft{}An open field; a meadow\textquotedblright{}.) Say it once more.
\end{tcolorbox}
\section[\texorpdfstring{\ja{しげる} (shigeru)}{しげる (shigeru)}]{\ja{しげる} (shigeru) — to grow thick}

\label{lesson:JA-C410-shigeru}



\begin{quote}
Before the new one: say the Japanese for a field (for crops), then the Japanese for a meadow.
\end{quote}

\subsection*{You'll want to know: \ja{しげる}}
\textbf{\ja{しげる}} — \emph{shigeru} — \textquotedblleft{}to grow thick\textquotedblright{}.

\begin{quote}
\textbf{\ja{なつは} \ja{くさが} \ja{しげります。}} — \emph{natsu wa kusa ga shigerimasu} — \textquotedblleft{}In summer the grass grows thick.\textquotedblright{}
\end{quote}

Polite \textbf{\ja{しげります}}, \emph{shigerimasu}; a godan verb. Past (polite): \textbf{\ja{にわの} \ja{くさが} \ja{しげりました。}}, \emph{niwa no kusa ga shigerimashita}.

\begin{cousinweb}
\textbf{\ja{しげる}} (\emph{shigeru}) is a native Japanese word (wago), inherited rather than borrowed from Chinese; its kanji was attached later as a kun'yomi (JMdict).
\end{cousinweb}

\begin{grammarlens}[title={What you've built}]
One more word for getting around.
\end{grammarlens}

\subsection*{Your turn}
\begin{itemize}
\raggedright
  \item \emph{Say it:} \emph{shigeru}
  \item \emph{Say it:} \emph{shigeru}, once more
  \item \emph{Say it:} \emph{natsu wa kusa ga shigerimasu} — In summer the grass grows thick.
  \item \emph{Recall:} say the Japanese for a field (for crops), then the Japanese for a meadow, then say \emph{shigeru} again
\end{itemize}

\begin{tcolorbox}[breakable,colback=teal!4,colframe=teal!35!black,title={Before you move on}]
What does \ja{しげる} mean? (\textquotedblleft{}To grow thick\textquotedblright{}.) Say it once more.
\end{tcolorbox}
